\vspace*{-1cm}
\section{Análisis del sistema}
\label{sec:analisis}


\subsection{Arquitectura del sistema}
Se elaborarán diagramas donde se indiquen las interfaces del sistema, los diferentes servicios que proporciona, los componentes en que se apoyan tales servicios, capa de acceso a datos, etc. Los diagramas de arquitectura NO deberían estar ligados a ninguna tecnología, sistema o implementación concreta.

\subsubsection{Arquitectura física}
Mostrará todos los componentes físicos que intervienen en un sistema para su despliegue y funcionamiento (servidores, conexiones, clientes…).

\subsubsection{Arquitectura lógica}
Si se trata de una arquitectura en tres capas, los diferentes subsistemas y módulos. Si se trata de una arquitectura basada en microservicios, cada microservicio, gestión de comunicaciones, gateway, etc.

\subsubsection{Diagrama de infraestructuras de nivel 2 (si procede)}
Se mostrarán los routers, switches, conexiones entre ellos, etc. Para cada uno de los dispositivos se indicarán sus características técnicas.

\begin{table}[h]
    \centering
    % chktex-disable 44
        \captionazul{Coste estimado del servicio}
        \setlength{\arrayrulewidth}{0.05mm} % <--- Grosor más fino (el estándar es 0.4pt)
        \begin{tabular}{|p{2cm}|p{10cm}|}
            \hline
            \textbf{Ref.} & \textbf{Descripción} \\
            \hline
            RF01 & Descripción del requisito funcional 01. \\
            \hline
            RF02 & Descripción del requisito funcional 02. \\
            \hline
            \ldots & \ldots \\
            \hline
            RFN & Descripción del requisito funcional N. \\
            \hline
        \end{tabular}
    \label{tab:coste_estimado_servicio}
\end{table}

\subsubsection{Diagrama de infraestructuras de nivel 3 (si procede)}
Se mostrarán las diferentes zonas de seguridad, firewalls, VLANs, servidores, etc. Los componentes de la infraestructura podrán ser tanto materiales como virtuales.


\subsection{Diseño de datos}
Diagrama Entidad-Relación extendido o equivalente. Se llevará un proceso de normalización y de optimización para obtener las tablas (modelo físico) de la base de datos.


\subsubsection{Migración y carga inicial de datos (si procede)}
En caso de que sea necesario, se elaborará un plan de migración y carga inicial de datos.

\subsection{Diseño de la interfaz de usuario}
Se mostrará cómo serán las ventanas, la navegación, etc. En caso de que sea necesario, podrán utilizarse diagramas de transición de estados.

\subsection{Diagrama de clases}
A lo largo de este documento se ha asumido un desarrollo orientado a objetos. No obstante, se pueden seguir otros enfoques, por ejemplo, un desarrollo estructurado con diagrama de contexto, diagramas de flujos de datos, diagrama entidad-relación, etc. Por otra parte, el tutor del proyecto puede dar sus propias orientaciones si lo considera oportuno.

\subsection{Entorno de construcción}
Se indicará el IDE, frameworks, simuladores, librerías, etc. que se han utilizado, así como otros detalles relevantes para la implementación del proyecto.

\subsection{Referencia al repositorio de software}
Se indicará la URL del repositorio en el que está el software para, si procede, que el tribunal pueda probarlo en sus propias máquinas. Si la aplicación puede ser desplegada en un sitio públicamente accesible, se deberá proporcionar la URL del sitio, así como toda la información necesaria para acceder al sistema.

Si este capítulo resultase demasiado extenso, podría separarse en dos.
